{\small
\renewcommand{\arraystretch}{1.15}
\begin{longtable}{@{}>{\raggedright\arraybackslash}p{0.22\dimexpr\linewidth-6\tabcolsep\relax}
                     >{\raggedright\arraybackslash}p{0.26\dimexpr\linewidth-6\tabcolsep\relax}
                     >{\raggedright\arraybackslash}p{0.22\dimexpr\linewidth-6\tabcolsep\relax}
                     >{\raggedright\arraybackslash}p{0.30\dimexpr\linewidth-6\tabcolsep\relax}@{}}
\caption{Parameter patterns in existing approaches}\label{tab:parameter-patterns}\\
\toprule
\textbf{Example approach} & \textbf{Representative information used} & \textbf{Role in model} & \textbf{Literature-derived implication} \\
\midrule
\endfirsthead
\toprule
\textbf{Example approach} & \textbf{Representative information used} & \textbf{Role in model} & \textbf{Literature-derived implication} \\
\midrule
\endhead
\bottomrule
\endlastfoot

Tailored centrality~\citep{simone2022,simone2023} &
Topology, intrinsic relevance, flow direction &
Improves structural ranking &
Topology can be enriched, but added variables need clear semantics \\
\addlinespace

Critical-pipe graph method~\citep{dastgir2023} &
Runoff/path information and capacity-related information &
Approximates criticality/flooding consequence &
Shows value of separating network paths from capacity response \\
\addlinespace

Weakest-link method~\citep{meijer2023} &
Storage, discharge capacity, flow paths &
Screens subsystem weakness &
Domain variables may be combined without redefining the generic graph concept \\
\addlinespace

Siltation-chain analysis~\citep{zhang2025} &
Pipeline risk, centrality, edge vulnerability, risk path &
Ranks critical nodes and paths &
Multiple parameters are justified when each represents a distinct mechanism \\
\addlinespace

GNN hydraulic surrogate~\citep{zhang2024} &
Hydraulic states, runoff, boundary, control variables, and graph topology &
Predicts physical hydraulic states &
High parameter count supports richer prediction but changes the task and data burden \\

\end{longtable}
}
